%% ma: L12-B02-0012
%% lop: 12
%% chuong: 1
%% bai: 2
%% dang: 12.2.8
%% loai: DS
%% muc: [TH, TH, VD, NB]
%% dap_an: "SSĐĐ"
%% nguon: "(NBV)-ĐỀ SỐ 19-MỖI NGÀY 1 ĐỀ THI 2025.pdf (D19, MỖI NGÀY 1 ĐỀ - Đề số 19), Phần II Câu 1 (Sở Ninh Bình 2025)"
%% kien_thuc: "Bài 2 (giá trị lớn nhất, nhỏ nhất của hàm số trên một đoạn), đạo hàm của hàm mũ, hàm căn thức"
%% trang_thai: cho-duyet
%% ly_do: "Dạng mới đề xuất 12.2.8 (hàm e^{u(x)} với u chứa căn thức: nghiệm f'(x)=0, công thức đạo hàm, giá trị lớn nhất và nhỏ nhất trên đoạn, giá trị tại đầu mút), chờ giáo viên duyệt."
%% ghi_chu: "Đáp án gốc S S Đ Đ khớp (giải độc lập bằng sympy: f'(x) = 0 chỉ có x = 2√2; max = e^{4√2}, min = e^{-4}; a + 2b + 3c = 10). Ý b): công thức đạo hàm đúng nhưng chỉ xác định trên khoảng (-4;4), không có tại x = ±4 nên 'với mọi x thuộc [-4;4]' là sai (đáp án gốc cũng chọn S, đề gốc nêu lí do này trong lời giải)."
\begin{ex}
Cho hàm số $y=f(x)=\mathrm e^{x+\sqrt{16-x^2}}$.
\choiceTF
{$f'(x)=0$ có hai nghiệm phân biệt.}
{$f'(x)=\left(1-\dfrac{x}{\sqrt{16-x^2}}\right)\mathrm e^{x+\sqrt{16-x^2}}$, $\forall x\in[-4;4]$.}
{\True Tích của giá trị lớn nhất và giá trị nhỏ nhất của hàm số $f(x)$ là $\mathrm e^{a+b\sqrt c}$ (với $a,b,c\in\mathbb Z$ và $c$ là số nguyên tố). Khi đó $a+2b+3c=10$.}
{\True $f(-4)=\dfrac1{\mathrm e^4}$.}
\loigiai{Hàm số xác định trên $[-4;4]$.

a) Với $x\in(-4;4)$: $f'(x)=\left(1-\dfrac{x}{\sqrt{16-x^2}}\right)\mathrm e^{x+\sqrt{16-x^2}}$. $f'(x)=0\Leftrightarrow\sqrt{16-x^2}=x\Leftrightarrow\begin{cases}x\ge0\\16-x^2=x^2\end{cases}\Leftrightarrow x=2\sqrt2$. Chỉ có một nghiệm nên a) sai.

b) Công thức trên chỉ đúng với $x\in(-4;4)$ vì hàm $\sqrt{16-x^2}$ không có đạo hàm tại $x=\pm4$. Do đó b) sai.

c) $f(2\sqrt2)=\mathrm e^{2\sqrt2+2\sqrt2}=\mathrm e^{4\sqrt2}$, $f(-4)=\mathrm e^{-4}$, $f(4)=\mathrm e^{4}$. Suy ra $\max f=\mathrm e^{4\sqrt2}$ ($4\sqrt2\approx5{,}66>4$) và $\min f=\mathrm e^{-4}$. Tích bằng $\mathrm e^{-4+4\sqrt2}$, tức $a=-4$, $b=4$, $c=2$ và $a+2b+3c=-4+8+6=10$. Vậy c) đúng.

d) $f(-4)=\mathrm e^{-4+\sqrt0}=\dfrac1{\mathrm e^4}$ nên d) đúng.}
\end{ex}
